\documentclass[10pt,a4paper]{article}
\usepackage[top=1.4cm, bottom=1.4cm, left=1.6cm, right=1.6cm, headheight=14pt, footskip=18pt]{geometry}
\usepackage[utf8]{inputenc}
\usepackage{amsmath,amsfonts,amssymb}
\usepackage{graphicx}
\usepackage{tikz}
\usepackage{circuitikz}
\usetikzlibrary{calc,patterns}
\usepackage[most]{tcolorbox}
\usepackage{enumitem}
\usepackage{fancyhdr}
\usepackage{booktabs}
\usepackage{tabularx}
\usepackage{colortbl}
\usepackage{hyperref}

% --- Palet Warna Resmi UNIB ---
\definecolor{unibblue}{RGB}{0, 32, 96}
\definecolor{unibgold}{RGB}{197, 160, 89}
\definecolor{unibgreen}{RGB}{34, 112, 60}
\definecolor{boxbg}{RGB}{248, 250, 253}
\definecolor{linegray}{RGB}{210, 215, 225}

% --- Pengaturan Header & Footer ---
\pagestyle{fancy}
\fancyhf{}
\lhead{\footnotesize\textbf{\color{unibblue}Universitas Bengkulu} \textbar\ Program Studi S1 Teknik Elektro}
\rhead{\footnotesize\color{gray}Persamaan Diferensial (Genap 2026)}
\lfoot{\footnotesize\color{gray}Problem Set Tugas Terstruktur Berbasis OBE}
\cfoot{\footnotesize\color{gray}Hal. \thepage\ dari 2}
\rfoot{\footnotesize\color{unibblue}\textbf{Minggu 7: Invers Transformasi Laplace \& Pe...}}
\renewcommand{\headrulewidth}{0.6pt}
\renewcommand{\footrulewidth}{0.4pt}

\begin{document}

% ==================== HALAMAN 1 ====================
\noindent\begin{minipage}[c]{0.68\textwidth}%
  {\large\textbf{\color{unibblue}PROBLEM SET (TUGAS MANDIRI TERSTRUKTUR)}}\\[1mm]
  {\normalsize\textbf{Mata Kuliah: Persamaan Diferensial (2 SKS)}}\\[0.5mm]
  {\small\textbf{Topik Minggu 7:} Invers Transformasi Laplace \& Pemutus Tenaga (TRV)}%
\end{minipage}%
\hfill%
\begin{minipage}[c]{0.30\textwidth}%
  \begin{tcolorbox}[colback=unibblue!5,colframe=unibblue,boxrule=0.6pt,arc=2pt,left=4pt,right=4pt,top=2pt,bottom=2pt]
    \scriptsize
    \textbf{Nama:} \dotfill\\[1.2mm]
    \textbf{NPM:} \dotfill\\[1.2mm]
    \textbf{Kelas/Klp:} \dotfill\\[1.2mm]
    \textbf{Nilai Final:} \hfill \textbf{/ 100}
  \end{tcolorbox}%
\end{minipage}\par

\vspace{1.5mm}

% Kotak Sub-CPMK & Pedoman Tugas Terstruktur
\begin{tcolorbox}[colback=boxbg,colframe=unibgold!90!black,boxrule=0.6pt,arc=2pt,left=5pt,right=5pt,top=3pt,bottom=3pt,title=\textbf{\scriptsize Capaian Pembelajaran (Sub-CPMK 4) \& Pedoman Pengerjaan Tugas}]
  \scriptsize
  \textbf{Sub-CPMK 4:} Mahasiswa mampu merekonstruksi sinyal waktu melalui Invers Transformasi Laplace (pecahan parsial \& integral konvolusi) serta menganalisis fenomena Transient Recovery Voltage (TRV) pada pemutus tenaga PMT.\\
  \textbf{Ketentuan Tugas:} Kerjakan secara mandiri pada kertas folio/A4 bergaris atau laporan digital rapi. Lampirkan penurunan rumus analitis lengkap dan cetak visualisasi kode program Julia (Soal C6). Kumpulkan pada awal perkuliahan minggu berikutnya.
\end{tcolorbox}

\vspace{1.5mm}

\noindent{\textbf{\color{unibblue}\large BAGIAN I: FONDASI KONSEPTUAL \& PENURUNAN MATEMATIS}}\par
\vspace{1mm}

% Soal C1
\noindent\textbf{1. (C1 - Mengingat \textbar\ Bobot: 10 Poin)}\par\vspace{0.8mm}
{\footnotesize Tuliskan metode pecahan parsial (*partial fraction expansion*) untuk fungsi rasional $F(s) = \frac{P(s)}{(s+a)(s^2 + \omega^2)}$, dan tuliskan pasangan invers transformasi Laplace untuk suku linier dan suku kuadratik tersebut.}\par

\vspace{2.5mm}

% Soal C2
\noindent\textbf{2. (C2 - Memahami \textbar\ Bobot: 15 Poin)}\par\vspace{0.8mm}
{\footnotesize Jelaskan definisi fisis *Transient Recovery Voltage* (TRV) pada Pemutus Tenaga (PMT / *Circuit Breaker*) menurut standar \textbf{IEC 62271-100}. Mengapa sesaat setelah arus busur listrik terputus pada titik arus nol (*current zero*), tegangan melintasi kontak PMT berosilasi dengan frekuensi tinggi?}\par

\vspace{2.5mm}

% Soal C3
\noindent\textbf{3. (C3 - Menerapkan \textbar\ Bobot: 20 Poin)}\par\vspace{0.8mm}
\noindent\begin{minipage}[t]{0.60\textwidth}%
  {\footnotesize Model domain-$s$ dari tegangan pemulihan transien PMT pada sistem transmisi 150 kV diberikan oleh $V_{\text{TRV}}(s) = V_m \left( \frac{1}{s} - \frac{s}{s^2 + \omega_0^2} \right)$. Gunakan invers transformasi Laplace untuk menurunkan fungsi waktu analitis $v_{\text{TRV}}(t)$, dan tentukan amplitudo puncak tegangan pemulihannya!}%
\end{minipage}%
\hfill%
\begin{minipage}[t]{0.37\textwidth}%
  \centering
  \resizebox{0.95\linewidth}{!}{%
  \begin{circuitikz}[american, scale=0.68, transform shape]
    \draw (0,0) to[sV, l=$v_s(t)$] (0,1.8)
          to[L, l=$L_{\text{trafo}}$] (1.6,1.8)
          to[switch, l={PMT Buka}] (3.0,1.8)
          to[short] (3.0,0)
          to[short] (0,0);
    \draw (1.6,1.8) to[C, l=$C_{\text{bus}}$, v=$v_{\text{TRV}}$] (1.6,0);
  \end{circuitikz}%
  }%
\end{minipage}\par

\vfill

% Kotak Formula Rujukan Teknis & Standar Industri
\begin{tcolorbox}[colback=unibblue!3,colframe=unibblue,colbacktitle=unibblue,coltitle=white,boxrule=0.6pt,arc=2pt,left=6pt,right=6pt,top=3pt,bottom=3pt,title=\textbf{\scriptsize FORMULA RUJUKAN TEKNIS \& STANDAR INDUSTRI TERKAIT}]
  \scriptsize
  \begin{itemize}[leftmargin=*, nosep]
  \item Ekspansi Pecahan Parsial: $\frac{P(s)}{(s-p_1)(s-p_2)} = \frac{A_1}{s-p_1} + \frac{A_2}{s-p_2}$ dengan $A_k = \lim_{s\to p_k}(s-p_k)F(s)$
  \item Teorema Konvolusi Domain Waktu: $\mathcal{L}^{-1}\{F(s)G(s)\} = (f * g)(t) = \int_0^t f(\tau) g(t-\tau) d\tau$
  \item Standar \textbf{IEC 62271-100}: Transient Recovery Voltage (TRV) pada Pemutus Tenaga (PMT)
  \item Laju Kenaikan Tegangan Pemulihan: $\text{RRRV} = \left. \frac{dv_{\text{TRV}}}{dt} \right|_{\text{maks}} = V_m \omega_0$
\end{itemize}
\end{tcolorbox}

\newpage
% ==================== HALAMAN 2 ====================

\noindent{\textbf{\color{unibblue}\large BAGIAN II: ANALISIS SISTEM, EVALUASI STANDAR, \& KOMPUTASI JULIA}}\par
\vspace{1mm}

% Soal C4
\noindent\textbf{4. (C4 - Menganalisis \textbar\ Bobot: 20 Poin)}\par\vspace{0.8mm}
{\footnotesize Analisis Laju Kenaikan Tegangan Pemulihan (\textit{Rate of Rise of Recovery Voltage} / RRRV) yang didefinisikan sebagai $\text{RRRV} = \left. \frac{dv_{\text{TRV}}}{dt} \right|_{\text{maks}}$. Jika induktansi sumber hubung singkat $L = 5\text{ mH}$, kapasitansi busbar $C = 20\text{ nF}$, dan tegangan puncak sistem $V_m = 120\text{ kV}$, hitung frekuensi osilasi alami $f_0$ dan nilai RRRV maksimum dalam satuan $\text{kV/}\mu\text{s}$.}\par

\vspace{3.5mm}

% Soal C5
\noindent\textbf{5. (C5 - Mengevaluasi \textbar\ Bobot: 20 Poin)}\par\vspace{0.8mm}
{\footnotesize Selesaikan Masalah Nilai Awal sistem eksitasi generator berikut secara tuntas menggunakan Transformasi Laplace: $y'' + 6y' + 25y = 50$, dengan nilai awal $y(0) = 1$ dan $y'(0) = -2$. Evaluasi apakah respons transien mengalami osilasi teredam, dan tentukan nilai tunak sistem saat $t \to \infty$.}\par

\vspace{3.5mm}

% Soal C6
\noindent\textbf{6. (C6 - Merancang \& Komputasi Julia \textbar\ Bobot: 15 Poin)}\par\vspace{0.8mm}
{\footnotesize Rancang skrip numerik berbasis \textbf{Julia} untuk menyelesaikan PDB TRV orde 2 non-homogen dengan paket \texttt{Differential\-Equations.jl} serta memplot kurva $v_{\text{TRV}}(t)$ bersama garis batas dielektrik PMT gas $\text{SF}_6$. Tuliskan definisi fungsi ODE dan parameter simulasi lengkap.}\par

\vfill

% Tabel Rekapitulasi Skor OBE & Lembar Catatan Umpan Balik Dosen
\noindent\begin{minipage}[b]{0.62\textwidth}%
  \centering
  \scriptsize
  \begin{tabular}{|c|c|c|c|c|c|c|}
    \hline
    \rowcolor{unibblue}
    \textbf{\color{white}C1} & \textbf{\color{white}C2} & \textbf{\color{white}C3} & \textbf{\color{white}C4} & \textbf{\color{white}C5} & \textbf{\color{white}C6} & \textbf{\color{white}Total} \\
    \rowcolor{unibblue!10}
    10 Pts & 15 Pts & 20 Pts & 20 Pts & 20 Pts & 15 Pts & 100 Pts \\
    \hline
    & & & & & & \\[3mm]
    \hline
  \end{tabular}%
\end{minipage}%
\hfill%
\begin{minipage}[b]{0.35\textwidth}%
  \centering
  \scriptsize
  \textbf{Verifikasi Dosen / Asisten:}\\[6mm]
  (\dotfill)\\[1mm]
  \tiny Tanggal: \dotfill
\end{minipage}\par

\vspace{2mm}
\begin{tcolorbox}[colback=white,colframe=unibblue!40!black,colbacktitle=unibblue!10,coltitle=unibblue,boxrule=0.5pt,arc=2pt,left=5pt,right=5pt,top=3pt,bottom=3pt,title=\textbf{\tiny Catatan Evaluasi \& Umpan Balik Dosen Pengampu}]
  \tiny
  \textbf{Rubrik Pemeriksaan Pengerjaan Tugas Mandiri:}\par\vspace{0.8mm}
  $\square$ Penurunan matematis langkah-demi-langkah runtut (C1--C4) \hfill $\square$ Validasi analitis \& standar industri IEEE/IEC/SPLN akurat (C5)\\[0.8mm]
  $\square$ Skrip program Julia mandiri \& grafik visualisasi terlampir (C6) \hfill $\square$ Kerapian format laporan \& ketepatan waktu pengumpulan\\[2.0mm]
  \textbf{Catatan / Koreksi Khusus Dosen:}\\[1.6cm]
\end{tcolorbox}

\end{document}
